\chapter{Abstract}
\label{ch:abstract}

Systems with interacting components are often modeled as networks, whether in the natural sciences, such as biology and physics, or in the social sciences. Of particular interest are transport processes within these systems, which are often studied using flow optimization techniques. This work focuses on flow in networks containing chemical reactions, referred to as \emph{reaction networks}, for which the molecular structures of the interacting species are known. These molecular structures influence the pathways that can be identified within the network. Reaction networks are naturally represented as hypergraphs, a mathematical formalism that facilitates the computational search for pathways.

The size, intricate connectivity, and combinatorial complexity of reaction networks make them difficult to analyze manually. A common query is to identify pathways leading to a specified target molecule from a given set of available source molecules while simultaneously satisfying the underlying physical and chemical constraints. Finding pathways that optimize the formation of a target molecule according to either thermodynamic favourability or reaction probability is a fundamental problem in many settings, including chemical reactor systems. This work makes three methodological contributions. First, it extends pathway search using thermodynamic information. Second, it incorporates kinetic information into pathway ranking. Third, it introduces a stochastic exploration framework for implicitly defined reaction networks.

This work integrates well-established computational chemistry tools and physically motivated modeling with existing algorithmic methods for querying pathways in reaction networks. It further extends the available executable graph transformation framework to enable a more realistic exploration of reaction networks. In particular, it presents a mixed-integer linear programming (MILP) formulation of the pathway search problem that incorporates chemical potentials and molecular concentrations. Because the molecular structures are known, chemical potentials can be estimated and, together with molecular concentrations, can be incorporated into the optimization model, allowing the search to be restricted to pathways consisting exclusively of thermodynamically favorable reactions. Furthermore, when multiple feasible pathways are identified, they can be ranked using a thermodynamically motivated objective function.

The computational methodology for identifying pathways in reaction networks is further extended by incorporating information about the kinetics of the individual reactions. The kinetic parameters are estimated from the molecular structures using a computational chemistry model. As multiple pathways frequently satisfy the specified search criteria, an objective function derived from physical principles is introduced to maximize the overall probability of a pathway. The resulting framework provides an automated pipeline for the flexible, kinetically informed investigation of pathways in large reaction networks, including networks generated through molecular space expansion methods that do not provide kinetic annotations.

Many reaction networks are too large to be generated and stored explicitly as molecular networks. To address this challenge, this work introduces a stochastic exploration framework for reaction networks that are implicitly defined by a graph grammar. The exploration strategy mimics stochastic chemical kinetics by combining collision-based pair selection with on-demand instantiation of reaction templates. At each exploration step, the algorithm first samples a pair of molecules to collide, then selects a reaction template, and finally samples one reaction instance from the set of matches of the selected template. This collision-first factorization avoids exhaustive enumeration of all currently possible reactions and enables efficient exploration of large reaction networks under both open- and closed-system conditions. The explored subnetwork can subsequently be regarded as a representative subset of the complete molecular space and analyzed using the methods developed in this work.

The proposed methods are demonstrated through several case studies that highlight both their strengths and their limitations. The thermodynamically informed framework is applied to a reaction network describing hydrogen cyanide–formamide chemistry. Alternative pathways to those currently accepted are identified and enumerated, including pathways that achieve more favorable objective values according to the thermodynamic ranking criterion. A reaction network generated from glycolonitrile, water, and ammonia is annotated with computed reaction energy barriers, after which pathways leading to glycine and glycolic acid are identified using the proposed methodology. Finally, the stochastic exploration framework is demonstrated on formose chemistry as a case study. The case study examines both the chemical characteristics of the explored subnetwork and the computational benefits of caching during exploration.
